\section{Equicontinuidad uniforme}

\begin{theorem}
  Sea \(E\) una familia de funciones derivables en un intervalo \(I\). Si existe \(c > 0\) tal que \(|f'(x)| \leq c\) para toda \(f \in E\) y todo \(x \in I\), entonces \(E\) es equicontinua en \(I\).
\end{theorem}

\begin{proof}
  Sean \(x_0 \in I\), \(\e > 0\) y \(\delta = \frac{\e}{c}\). Si \(x \in I\) y \(|x - x_0| < \delta\), por el teorema del valor medio, para toda \(f \in E\),
  \[
    |f(x) - f(x_0)| \leq c \, |x - x_0| < c \cdot \frac{\e}{c} = \e. \qedhere
  \]
\end{proof}

\begin{note}
  Con el mismo argumento, si para todo \(x \in I\) existen \(c_x > 0\) y un intervalo \(I_x \subseteq I\) que contiene a \(x\) en su interior relativo a \(I\), tales que \(|f'(y)| \leq c_x\) para toda \(f \in E\) y todo \(y \in I_x\), entonces \(E\) es equicontinua en \(I\). Basta tomar \(\delta = \min\{ \frac{\e}{c_x}, \eta \}\), con \(\eta > 0\) tal que \(I \cap (x - \eta, x + \eta) \subseteq I_x\).
\end{note}

\begin{definition}[Equicontinuidad uniforme]
  Una familia \(E\) de funciones \(f\colon X \to \R\) se dice \emph{uniformemente equicontinua} si para todo \(\e > 0\) existe \(\delta > 0\) tal que, si \(x, y \in X\) y \(|x - y| < \delta\), entonces \(|f(x) - f(y)| < \e\) para toda \(f \in E\).
\end{definition}

\begin{eg}
  Si \(E\) es una familia de funciones derivables en un intervalo \(I\) con \(|f'(x)| \leq c\) para toda \(f \in E\) y todo \(x \in I\), entonces \(E\) es uniformemente equicontinua: el mismo argumento con \(\delta = \frac{\e}{c}\) vale para todo par \(x, y \in I\).
\end{eg}

\begin{eg}
  Una familia formada por una única función continua que no es uniformemente continua, por ejemplo \(f(x) = x^2\) en \(\R\), es equicontinua pero no uniformemente equicontinua.
\end{eg}

\begin{theorem}
  Sea \(K \subset \R\) compacto. Toda familia equicontinua de funciones \(f\colon K \to \R\) es uniformemente equicontinua.
\end{theorem}

\begin{proof}
  Supongamos que una familia equicontinua \(E\) no es uniformemente equicontinua. Entonces existe \(\e > 0\) tal que para todo \(n \in \N\) hay \(x_n, y_n \in K\) y \(f_n \in E\) con \(|x_n - y_n| < \frac{1}{n}\) y \(|f_n(x_n) - f_n(y_n)| \geq \e\).

  Como \(K\) es compacto, pasando a una subsucesión podemos suponer que \(x_n \to x \in K\), y como \(|y_n - x_n| < \frac{1}{n}\), también \(y_n \to x\). Por la equicontinuidad de \(E\) en \(x\) existe \(\delta > 0\) tal que \(z \in K\) y \(|z - x| < \delta\) implican \(|f(z) - f(x)| < \frac{\e}{2}\) para toda \(f \in E\). Para \(n\) suficientemente grande, \(|x_n - x| < \delta\) y \(|y_n - x| < \delta\), de modo que
  \[
    |f_n(x_n) - f_n(y_n)| \leq |f_n(x_n) - f_n(x)| + |f_n(x) - f_n(y_n)| < \e,
  \]
  lo que es absurdo.
\end{proof}

\begin{theorem}
  Sea \({(f_n)}\) una sucesión equicontinua de funciones \(f_n\colon X \to \R\), con \(X \subseteq \R\). Si \({(f_n)}\) converge puntualmente en un subconjunto denso \(D \subseteq X\), entonces converge uniformemente sobre todo compacto \(K \subseteq X\).
\end{theorem}

\begin{proof}
  Sean \(K \subseteq X\) compacto y \(\e > 0\). Veamos que existe \(n_0\) tal que \(|f_m(x) - f_n(x)| < \e\) para todo \(m, n > n_0\) y todo \(x \in K\).

  Para cada \(y \in K\), por la equicontinuidad en \(y\) existe un intervalo abierto \(I_y\) centrado en \(y\) tal que \(|f_n(z) - f_n(y)| < \frac{\e}{6}\) para todo \(z \in X \cap I_y\) y todo \(n \in \N\); luego \(|f_n(x) - f_n(z)| < \frac{\e}{3}\) para todo \(x, z \in X \cap I_y\) y todo \(n\). Como \(K\) es compacto, del cubrimiento \(K \subseteq \bigcup_{y \in K} I_y\) se puede extraer un subcubrimiento finito \(K \subseteq I_1 \cup \cdots \cup I_p\), donde cada \(I_i = I_{y_i}\) con \(y_i \in K \subseteq X\).

  Como \(D\) es denso en \(X\) y \(y_i \in X \cap I_i\), que es abierto relativo a \(X\), podemos elegir \(d_i \in D \cap I_i\) para cada \(i\). Como \(f_n(d_i)\) converge, existe \(n_i\) tal que \(|f_m(d_i) - f_n(d_i)| < \frac{\e}{3}\) para todo \(m, n > n_i\). Sea \(n_0 = \max\{n_1, \dots, n_p\}\).

  Si \(m, n > n_0\) y \(x \in K\), existe \(i\) con \(x \in I_i\), y entonces
  \[
    |f_m(x) - f_n(x)| \leq |f_m(x) - f_m(d_i)| + |f_m(d_i) - f_n(d_i)| + |f_n(d_i) - f_n(x)| < \e.
  \]
  Luego \({(f_n)}\) es uniformemente de Cauchy en \(K\) y, por lo tanto, converge uniformemente en \(K\).
\end{proof}

\begin{definition}[Acotación puntual y uniforme]
  Una familia \(E\) de funciones \(f\colon X \to \R\) se dice \emph{puntualmente acotada} si para todo \(x \in X\) existe \(c_x > 0\) tal que \(|f(x)| \leq c_x\) para toda \(f \in E\). Se dice \emph{uniformemente acotada} si existe \(c > 0\) tal que \(|f(x)| \leq c\) para toda \(f \in E\) y todo \(x \in X\).
\end{definition}

\begin{theorem}[Cantor, Tijonov]
  Sea \(X \subseteq \R\) numerable. Toda sucesión puntualmente acotada de funciones \(f_n\colon X \to \R\) posee una subsucesión puntualmente convergente.
\end{theorem}

\begin{proof}
  Si \(X\) es finito el resultado se sigue de que toda sucesión acotada en \(\R\) tiene una subsucesión convergente, aplicada sucesivamente a cada punto. Sea \(X = \{x_1, x_2, \dots\}\) infinito. Como \({(f_n(x_1))}_{n \in \N}\) es acotada, tiene una subsucesión convergente: existe \(\N_1 \subseteq \N\) infinito tal que existe \(a_1 = \lim_{n \in \N_1} f_n(x_1)\). Como \({(f_n(x_2))}_{n \in \N_1}\) es acotada, existe \(\N_2 \subseteq \N_1\) infinito tal que existe \(a_2 = \lim_{n \in \N_2} f_n(x_2)\). Siguiendo así, para cada \(i \in \N\) obtenemos \(\N_1 \supseteq \N_2 \supseteq \cdots \supseteq \N_i \supseteq \cdots\), infinitos, con \(\lim_{n \in \N_i} f_n(x_i) = a_i\).

  Definimos \(\N^* \subseteq \N\) tomando como \(i\)-ésimo elemento de \(\N^*\) al \(i\)-ésimo elemento de \(\N_i\) (en orden creciente). Como \(\N_{i+1} \subseteq \N_i\), el \((i+1)\)-ésimo elemento de \(\N_{i+1}\) es mayor o igual que el \((i+1)\)-ésimo de \(\N_i\), que es mayor que el \(i\)-ésimo de \(\N_i\); luego \(\N^*\) es infinito y estrictamente creciente. Para cada \(i\), a partir del \(i\)-ésimo elemento, \({(f_n(x_i))}_{n \in \N^*}\) es una subsucesión de \({(f_n(x_i))}_{n \in \N_i}\), que converge. Por lo tanto \({(f_n)}_{n \in \N^*}\) converge en todo \(x_i \in X\).
\end{proof}

\begin{theorem}[Arzelà, Ascoli]
  Sea \(K \subset \R\) compacto. Toda sucesión equicontinua y puntualmente acotada de funciones \(f_n\colon K \to \R\) posee una subsucesión uniformemente convergente.
\end{theorem}

\begin{proof}
  Veamos primero que \(K\) tiene un subconjunto numerable y denso \(X\). Para cada \(m \in \N\), los intervalos abiertos de radio \(\frac{1}{m}\) centrados en puntos de \(K\) cubren \(K\); por compacidad bastan finitos, con centros en un conjunto finito \(X_m \subseteq K\). Entonces \(X = \bigcup_m X_m\) es numerable y todo punto de \(K\) está a distancia menor que \(\frac{1}{m}\) de \(X_m\) para cada \(m\), así que \(X\) es denso en \(K\).

  Por el teorema de Cantor--Tijonov, como \(X\) es numerable y \({(f_n)}\) es puntualmente acotada, existe una subsucesión \({(f_{n_k})}_{k \in \N}\) que converge puntualmente en \(X\). Esta subsucesión sigue siendo equicontinua, y converge puntualmente en el subconjunto denso \(X \subseteq K\). Por el teorema anterior, converge uniformemente sobre el compacto \(K\).
\end{proof}

\section{Ecuaciones diferenciales}

TODO
